\documentclass[11pt]{article}
\usepackage[margin=1in]{geometry}
\usepackage{graphicx}
\usepackage{booktabs}
\usepackage{amsmath}
\usepackage{amssymb}
\usepackage{natbib}
\usepackage{hyperref}
\usepackage{xcolor}
\usepackage{times}
\hypersetup{colorlinks=true, linkcolor=blue, citecolor=blue, urlcolor=blue}

\title{Systemic pharmacological suppression of neural activity reverses
oculomotor learning impairment in a Purkinje-cell-specific \textit{Fmr1}
knockout model of Fragile X syndrome}
\author{}
\date{}

\begin{document}
\maketitle

\begin{abstract}
Enhancement of associative synaptic plasticity frequently impairs rather than
enhances learning, a paradox that obstructs the clinical translation of
plasticity-enhancing interventions for learning disorders. The threshold
metaplasticity hypothesis, developed in mice lacking the major
histocompatibility class I molecules H2-K\textsuperscript{b} and H2-D\textsuperscript{b}
(MHCI KbDb\textsuperscript{-/-}), proposes that enhanced long-term depression
(LTD) at parallel fibre--Purkinje cell (PF--Purkinje cell) synapses is
aberrantly recruited by ongoing circuit activity, raising the threshold for
its further induction and reducing its availability to support learning.
Here we tested this hypothesis in a second, molecularly distinct genetic
model with enhanced PF--Purkinje cell LTD --- Purkinje-cell-specific
\textit{Fmr1} knockout mice (L7-Fmr1 KO) --- and asked whether a systemic
pharmacological manipulation could reset the circuit. L7-Fmr1 KO mice were
selectively impaired on high-frequency (1 Hz) vestibulo-ocular reflex
(VOR) increase learning ($p = 0.97$, 0 vs. 30 min; $p = 3.47 \times 10^{-5}$
vs. wild-type) and optokinetic reflex (OKR) adaptation
($p = 1.27 \times 10^{-4}$ vs. wild-type at 60 min), while VOR-decrease
learning was intact ($p = 1.10 \times 10^{-5}$, 0 vs. 30 min). Behavioural
pre-training designed to reverse LTD rescued VOR-increase learning
($p = 0.001$ after VOR-decrease pre-training), and a single intraperitoneal
injection of diazepam (0.5 mg/kg) 18--24~h before training fully rescued
both VOR-increase ($p = 0.86$ vs. wild-type) and OKR adaptation
($p = 0.02$, with KO learning exceeding WT); the rescue was transient and
task-specific, leaving low-frequency (0.5 Hz) VOR-increase learning
($p = 0.66$) and VOR-decrease learning unchanged. Diazepam did not alter
baseline VOR or OKR gain and did not enhance learning in wild-type
littermates. These results extend the threshold metaplasticity hypothesis
to a molecularly independent model, support a generalisable pharmacological
strategy for resetting circuits with pathologically enhanced associative
LTD, and motivate the repurposing of approved GABAergic drugs for
learning disorders associated with enhanced plasticity.
\end{abstract}

\section{Introduction}

Long-term synaptic plasticity has long been pursued as a therapeutic target
for disorders of learning and memory, and a broad toolkit now exists for
augmenting it. Yet the relationship between enhanced plasticity and improved
behaviour has proven deeply ambiguous: in some studies, enhancing plasticity
improves learning, while in many others the same class of manipulations
produces learning deficits \citep{koekkoek2005fmr1, nosyreva2006fmr1,
nguyenvu2017saturation}. This ambiguity is not an idle curiosity. It has
direct consequences for the clinical translation of plasticity-enhancing
interventions for stroke recovery, neurodegenerative disease, addiction,
and developmental disorders such as Fragile X syndrome (FXS). A mechanistic
account of when and why augmented plasticity fails to support adaptive
learning would constrain how such interventions are designed and deployed.

A particularly informative vantage point on this puzzle is the cerebellum,
where associative long-term depression (LTD) at parallel-fibre--Purkinje
cell (PF--Purkinje cell) synapses has long been implicated in motor learning
\citep{boyden2004cerebellum, raymond2018cerebellum}. The cerebellum supports
a family of closely related oculomotor learning tasks that share vestibular,
visual, and motor pathways but differ in their reliance on PF--Purkinje cell
LTD. Increases in the amplitude of the vestibulo-ocular reflex (VOR-increase
learning) and optokinetic reflex (OKR) adaptation appear to require
PF--Purkinje cell LTD
\citep{boyden2004cerebellum, inoshita2018okr, takeuchi2008delphilin,
kimpo2014gating, rowan2018graded}, whereas VOR-decrease learning does not
\citep{boyden2004cerebellum, raymond2018cerebellum, kimpo2005vor,
kimpo2014gating}. This dissociation provides a within-subject internal
control that is unusually stringent for a plasticity--learning link.

The paradox of enhanced plasticity emerges sharply in this system. Genetic
manipulations that abolish PF--Purkinje cell LTD, as well as those that
enhance it, both impair the same subset of oculomotor learning tasks
\citep{nguyenvu2017saturation, koekkoek2005fmr1}. A satisfying account of
this symmetry must explain how more plasticity and less plasticity produce
the same behavioural phenotype. \citet{nguyenvu2017saturation} proposed
such an account for mice lacking the two major histocompatibility class I
molecules H2-K\textsuperscript{b} and H2-D\textsuperscript{b}
(MHCI KbDb\textsuperscript{-/-}), which exhibit enhanced PF--Purkinje cell
LTD. They argued that when LTD is lowered in threshold, spontaneous activity
in the cerebellar circuit before training aberrantly recruits the plasticity
machinery, saturating or, more generally, raising the threshold for its
subsequent induction. During training, LTD is then effectively unavailable
at the synapses that need it to support learning. Two predictions of this
\emph{threshold metaplasticity} hypothesis were confirmed behaviourally:
optogenetic induction of LTD in wild-type (WT) animals recapitulates the
oculomotor learning deficit, and behavioural pre-training designed to
reverse LTD in vivo rescues learning in MHCI KbDb\textsuperscript{-/-} mice
\citep{nguyenvu2017saturation}.

The current study extends this logic along two axes. First, we test the
threshold metaplasticity hypothesis in a second mouse line with enhanced
PF--Purkinje cell LTD arising from an entirely different molecular pathway.
Purkinje-cell-specific deletion of the Fragile X gene \textit{Fmr1} enhances
PF--Purkinje cell LTD and is a well-characterised model of the cerebellar
component of FXS \citep{koekkoek2005fmr1, mientjes2006conditional,
zhang2004l7cre, nosyreva2006fmr1}. Whereas MHCI KbDb\textsuperscript{-/-}
animals act on the MAP-kinase and integrin pathways, L7-\textit{Fmr1}
knockout (L7-Fmr1 KO) mice act via loss of FMRP-mediated translational
control at dendritic synapses \citep{koekkoek2005fmr1, nosyreva2006fmr1}.
A shared learning phenotype across molecularly distinct lines would strongly
implicate their common downstream effect --- enhanced LTD --- as the cause
of impairment, rather than a coincident off-target perturbation. Second,
we introduce a new pharmacological prediction that was not tested in the
prior work: because PF--Purkinje cell LTD requires coincident parallel-fibre
and climbing-fibre activity \citep{levram2003reversing}, a systemic drug
that suppresses neural activity for a sustained window before training
should, upon recovery, restore learning capacity. We used the
FDA-approved positive allosteric GABA\textsubscript{A} modulator diazepam
for this purpose and assessed learning after the drug had cleared
(18--24~h post-injection), separating any transient sedative effect from
lasting changes in the metaplastic state of the circuit.

Our experimental design is embedded within the sliding-threshold framework
formulated by Bienenstock, Cooper, and Munro (BCM) and refined in
subsequent theoretical and experimental studies of metaplasticity
\citep{abraham1996metaplasticity}. That framework was built around Hebbian
LTP, for which a threshold that rises with activity provides a homeostatic
counterweight to runaway potentiation. Associative LTD at PF--Purkinje
cell synapses is anti-Hebbian, and a similar sliding-threshold mechanism
would not serve firing-rate homeostasis. Our results therefore also bear
on the conceptual question of what threshold metaplasticity does when the
underlying plasticity is not destabilising in the BCM sense.

We report three principal findings. First, L7-Fmr1 KO mice are selectively
impaired on high-frequency VOR-increase learning and on OKR adaptation,
with intact VOR-decrease learning --- mirroring the pattern previously
described in MHCI KbDb\textsuperscript{-/-} mice. Second, behavioural
pre-training fully rescues the high-frequency VOR-increase deficit, as
predicted by threshold metaplasticity. Third, a single systemic dose of
diazepam, given 18--24~h before training, achieves the same rescue for both
VOR-increase and OKR adaptation, exceeds WT performance on OKR, but has no
effect on WT learning, on VOR-decrease learning in either genotype, or on
low-frequency VOR-increase learning. These results establish pre-training
activity suppression as a general, pharmacologically tractable strategy
for resetting circuits with pathologically enhanced associative LTD, and
motivate its exploration in other neurological and psychiatric conditions.

\section{Related work}

\paragraph{PF--Purkinje cell LTD and selective cerebellar learning.}
LTD at the PF--Purkinje cell synapse has been advanced as a central plasticity
mechanism for cerebellar motor learning for decades
\citep{boyden2004cerebellum}. Contemporary evidence, however, points to a
selective rather than universal role: PF--Purkinje cell LTD is strongly
implicated in VOR-increase learning with high-frequency ($\geq 1$~Hz)
stimuli and in OKR adaptation, while VOR-decrease learning and low-frequency
($\leq 0.66$~Hz) VOR-increase learning proceed normally in LTD-impaired
animals \citep{boyden2004cerebellum, raymond2018cerebellum, kimpo2005vor,
kimpo2014gating, inoshita2018okr, takeuchi2008delphilin, rowan2018graded,
nguyenvu2017saturation}. This task specificity provides a sharp behavioural
fingerprint of LTD dependence that we exploit in the present experiments.

\paragraph{Enhanced PF--Purkinje cell LTD in Fragile X and MHCI models.}
Two genetic models produce enhanced PF--Purkinje cell LTD: the
MHCI KbDb\textsuperscript{-/-} double knockout and Purkinje-cell-specific
\textit{Fmr1} knockout \citep{koekkoek2005fmr1, nosyreva2006fmr1,
nguyenvu2017saturation}. Both exhibit enlarged dendritic spines and atypical
synaptic signatures, and both yield oculomotor learning impairments
selectively on the tasks most strongly linked to LTD
\citep{nguyenvu2017saturation, koekkoek2005fmr1}. The Fragile X model
carries the additional clinical motivation that patients with FXS display
cerebellar contributions to eyeblink conditioning deficits and gait ataxia
\citep{koekkoek2005fmr1}, making translational strategies that rescue
cerebellar learning directly relevant.

\paragraph{Reversal of LTD and behavioural pre-training.}
PF--Purkinje cell LTD can be experimentally reversed by the induction of
PF--Purkinje cell LTP, which is mediated by nitric oxide and postsynaptic
calcium elevation \citep{levram2003reversing}. In the intact cerebellum,
VOR-decrease training is thought to engage PF--Purkinje cell LTP and has
been used as a behavioural pre-training intervention to reset the plasticity
state of the circuit. \citet{nguyenvu2017saturation} showed that a session
of VOR-decrease pre-training rescues VOR-increase learning in
MHCI KbDb\textsuperscript{-/-} mice, without affecting WT performance.
Whether this behavioural reset generalises to a molecularly distinct
LTD-enhancing manipulation has not been directly tested, and is one
objective of the present work.

\paragraph{Metaplasticity and the BCM framework.}
The concept of an experience-dependent threshold for synaptic plasticity
has been central to theoretical neuroscience since the BCM model, which
posited that the modification threshold for LTP slides with recent activity
to stabilise circuits against runaway potentiation
\citep{abraham1996metaplasticity}. A large body of experimental work on
NMDA-receptor-dependent Hebbian plasticity in cortex and hippocampus is
broadly consistent with this picture, but direct evidence for threshold
metaplasticity in the context of anti-Hebbian, associative LTD has been
limited. The molecular substrates that could implement such a threshold at
the PF--Purkinje cell synapse are, however, plentiful and well-characterised:
plasticity of the climbing-fibre synapse, Purkinje-cell dendritic
excitability mediated by SK2 channels \citep{ohtsuki2012sk2}, modulation
of climbing-fibre burst structure, and activity-dependent plasticity of
molecular-layer interneuron inhibition onto Purkinje cells
\citep{rowan2018graded}. PF--Purkinje cell LTD itself is mediated by
AMPA-receptor trafficking via stargazin (TARP) dephosphorylation
\citep{nomura2012stargazin}.

\paragraph{Memory allocation and threshold metaplasticity.}
Threshold metaplasticity is closely related to theoretical work on memory
allocation, which asks how particular memories become preferentially encoded
by particular neurons. In the lateral amygdala, CREB-dependent increases in
dendritic spine density bias neurons toward inclusion in a fear-memory trace
\citep{sargin2013creb}. A complementary proposal is that a plasticity-driven
\emph{increase} in threshold can \emph{exclude} synapses, and therefore
neurons, from subsequent memory traces, thereby protecting recently acquired
memories and separating temporally proximate ones. The BCM framework, with
its explicit link between activity history and plasticity threshold,
provides a natural substrate for such allocation rules
\citep{abraham1996metaplasticity}.

\section{Methods}

\subsection{Animals}

Purkinje-cell-specific \textit{Fmr1} knockout mice were generated by
crossing homozygous floxed \textit{Fmr1} females (\textit{Fmr1} conditional
knockout allele; \citealp{mientjes2006conditional}) with males expressing
L7/Pcp2-Cre on an autosome \citep{zhang2004l7cre}. L7/Pcp2-Cre drives
Cre recombinase expression in a manner that is highly selective for
cerebellar Purkinje cells. Male offspring from the initial cross were then
mated with \textit{Fmr1}-floxed homozygous females to generate offspring
homozygous for the floxed \textit{Fmr1} allele, half of which were
Cre-positive. Cre-positive offspring (L7-Fmr1 KO) and their Cre-negative
littermates (WT) were used throughout. Genotyping was performed by
Transnetyx Inc. on ear-clip samples using RT-qPCR for the floxed
\textit{Fmr1} allele and for Cre. Mice were housed on a reversed 12-h
light/dark cycle with behavioural testing during the dark phase. Male and
female mice 8--22 weeks old were used; learning deficits were similar
across sexes and data were pooled across sex.

\subsection{Surgery and eye-movement recording}

Mice underwent surgery between 8--12 weeks of age to implant head-restraint
and eye-movement recording hardware, as described previously
\citep{nguyenvu2017saturation}. Under 1.5--2.5\% isoflurane anaesthesia, a
scalp incision was made and a custom head post was attached to the skull
with dental acrylic. Two stacked neodymium magnets (total $0.75 \times 2$~mm,
grade N50, axially magnetised) were implanted beneath the conjunctiva on
the temporal side of the left eye, and an angular magnetic field sensor
(Honeywell HMC1512) was soldered to an 8-pin connector and attached to
the skull above the eye in a plane parallel to horizontal (nasal--temporal)
eye movements. Eye movements were measured by detecting changes in the
magnetic field caused by the implanted magnets moving with the eye. Mice
recovered for at least five days before experiments began.

\subsection{Oculomotor tasks and stimuli}

Experiments were performed in a light-proof, sound-attenuated chamber with
the head restrained on a vestibular turntable. Vestibular stimuli were
sinusoidal rotations of the turntable about an earth-vertical axis at
1~Hz or 0.5~Hz with a peak velocity of $\pm 10^{\circ}/\text{s}$. Visual
stimuli were provided by an optokinetic drum with alternating black and
white stripes ($\sim 7.5^{\circ}$ each) illuminated by an LED strip,
rotated independently about the same axis.

\paragraph{VOR testing and training blocks.}
A testing block consisted of three 45-s tests of the VOR, measured in
complete darkness with the turntable in sinusoidal rotation; the three
tests were separated by 10-s inter-trial intervals. A training block was
10 min long. For VOR-increase training, the vestibular turntable rotation
was paired with oppositely directed motion of the illuminated optokinetic
drum at the same peak velocity. For VOR-decrease training, the drum was
moved in the same direction as the turntable with the same velocity,
rendering the drum stationary relative to the head. In pre-training
experiments, the pre-training consisted of three 10-min blocks of either
VOR-decrease training or vestibular-only stimulation (turntable rotation
alone in the dark), each preceded by a testing block. Training totalled
30 min across three 10-min blocks interleaved with testing blocks.

\paragraph{OKR adaptation.}
OKR adaptation consisted of sixty 50-s blocks of 1~Hz sinusoidal optokinetic
drum rotation at $\pm 10^{\circ}/\text{s}$ with the turntable stationary,
each block followed by 10~s in darkness.

\subsection{Drug administration}

In drug experiments, mice received a single intraperitoneal (IP) injection
of 0.4, 0.5, or 2.5 mg/kg diazepam (dissolved in saline), or a saline
vehicle control. After injection, mice returned to the home cage and
behavioural testing was performed 2~h, 18--24~h, or 1 week later. The
order of drug and saline sessions was pseudorandomised, with at least two
days between consecutive experiments. The 0.5 mg/kg dose with the
18--24-h delay was used for all rescue experiments reported in the Results.

\subsection{Eye-movement analysis}

Eye-position signals from the magnetic sensor were filtered with a
fourth-order 15~Hz Butterworth low-pass filter and digitally differentiated
with a 30-ms windowed Savitzky--Golay filter to obtain eye velocity. Eye
velocity within each 45-s VOR test or 50-s OKR block was fit with a
sinusoid at the stimulus frequency (1~Hz or 0.5~Hz). Points deviating from
the sinusoidal fit by more than $31^{\circ}/\text{s}$ were flagged as
saccades or artefacts and excluded together with the 50~ms before and after;
segments shorter than 10~ms were also excluded. If more than 45\% of an
entire VOR test or OKR block was excluded, that test/block was discarded
in full. The remaining data were refit with a sinusoid and the amplitude
used as the measure of the eye-movement response. The VOR amplitude for
each training block was the average over the three 45-s tests in the
following testing block. VOR learning ($\Delta$VOR) was expressed as the
percentage change in VOR amplitude from the pre-training baseline. For
OKR, the pre-training baseline was the mean of the first three 50-s
blocks, and the post-training measure was the mean of the last three
blocks (blocks 58--60); OKR learning ($\Delta$OKR) was the percentage
change from pre- to post-training. Eye-movement gain was defined as the
ratio of eye-movement amplitude to stimulus amplitude for the VOR or the
OKR, respectively.

\subsection{Statistical analysis}

Normality was tested using the Shapiro--Wilk test. Group comparisons were
performed with a two-factor repeated-measures ANOVA followed by Tukey
post-hoc tests, or with two-sample or paired-sample $t$-tests as
appropriate. All analyses were performed in OriginPro 2022. A threshold
of $p < 0.05$ was used throughout. Data are reported as mean $\pm$~SEM
unless noted otherwise. All experimental procedures were approved by the
Administrative Panel on Laboratory Animal Care at Stanford University.

\section{Results}

\subsection{L7-Fmr1 KO mice are selectively impaired on VOR-increase
learning and OKR adaptation}

To determine whether enhanced PF--Purkinje cell LTD, produced by a
Purkinje-cell-restricted loss of \textit{Fmr1}, recapitulates the
task-selective oculomotor learning deficits described in
MHCI KbDb\textsuperscript{-/-} animals, we compared L7-Fmr1 KO mice to
Cre-negative WT littermates on two VOR paradigms that are differentially
sensitive to PF--Purkinje cell LTD \citep{boyden2004cerebellum,
raymond2018cerebellum}. Mice were head-restrained on a vestibular turntable
and eye velocity was recorded magnetically during 30 min of training in
which a 1~Hz, $\pm 10^{\circ}/\text{s}$ vestibular rotation was paired with
either oppositely-directed (VOR-increase) or same-directed (VOR-decrease)
motion of a large-field optokinetic drum.

WT mice acquired both forms of learning over 30 min of training. VOR-increase
training produced a significant increase in VOR amplitude ($p = 7.37 \times
10^{-4}$, 0 vs. 30 min, Tukey post-hoc), and VOR-decrease training produced
a significant decrease ($p = 0.001$). L7-Fmr1 KO animals showed a striking
and selective deficit: they exhibited no change in VOR amplitude after 30
min of VOR-increase training ($p = 0.97$, 0 vs. 30 min in KO;
$p = 3.47 \times 10^{-5}$ for KO vs. WT at 30 min), but they showed
completely normal VOR-decrease learning ($p = 1.10 \times 10^{-5}$, 0 vs.
30 min in KO; $p = 0.091$ for KO vs. WT at 30 min) (Figure~\ref{fig:selective_impairment}).
The VOR-increase deficit was present in both sexes (male $p = 0.01$;
female $p = 0.05$; KO vs. WT at 30 min), so data were pooled across sex
throughout.

Baseline oculomotor performance was normal. VOR gain in the dark before
training did not differ between genotypes (two-sample $t$-test, $p = 0.95$),
and the eye-movement responses to the paired visual--vestibular stimuli
during the first 45~s of each training type were also genotype-matched
(VOR-increase paired stimuli, $p = 0.50$; VOR-decrease paired stimuli,
$p = 0.76$). These controls rule out vestibular, visual, or oculomotor
performance confounds and localise the L7-Fmr1 KO phenotype to a selective
learning deficit on VOR-increase training. The pattern --- intact
VOR-decrease alongside profoundly impaired VOR-increase --- mirrors the
oculomotor phenotype of MHCI KbDb\textsuperscript{-/-} mice
\citep{nguyenvu2017saturation}, providing the first test that a
molecularly distinct enhancement of PF--Purkinje cell LTD yields the same
selective behavioural signature.

\begin{figure}[htbp]
\centering
\includegraphics[width=0.95\textwidth]{figures/fig1_selective_impairment.pdf}
\caption{\textbf{Selective VOR-increase learning impairment in L7-Fmr1 KO mice.}
Bars show $-\log_{10}$ of the Tukey post-hoc $p$-values extracted from the
VOR experiments with 1~Hz sinusoidal vestibular--visual pairing. Green
bars indicate $p<0.05$ (significant difference), red bars indicate
not-significant. WT mice learned both VOR-increase
($p=7.37\times10^{-4}$, 0 vs.\ 30 min) and VOR-decrease ($p=0.001$) tasks;
L7-Fmr1 KO mice failed on VOR-increase ($p=0.97$; KO vs.\ WT
$p=3.47\times10^{-5}$) but learned VOR-decrease normally
($p=1.10\times10^{-5}$; KO vs.\ WT $p=0.091$).}
\label{fig:selective_impairment}
\end{figure}

\begin{figure}[htbp]
\centering
\includegraphics[width=0.95\textwidth]{figures/fig2_behavioral_pretraining.pdf}
\caption{\textbf{Behavioural pre-training rescues VOR-increase learning in
L7-Fmr1 KO mice.} Extracted Tukey or paired-$t$-test $p$-values. Without
pre-training, KO animals did not learn ($p=0.99$, 0 vs.\ 30 min); after
30 min of VOR-decrease pre-training, they did ($p=0.001$). Vestibular-only
pre-training produced a KO vs.\ WT comparison of $p=0.84$. The within-mouse
comparison (no-pre vs.\ pre) was $p=0.01$ in the KO cohort. Subsampling for
matched pre-training-induced decreases preserved the rescue.}
\label{fig:behavioral_pretraining}
\end{figure}

\begin{figure}[htbp]
\centering
\includegraphics[width=0.95\textwidth]{figures/fig3_diazepam_rescue.pdf}
\caption{\textbf{Diazepam pre-treatment (0.5 mg/kg IP, 18--24~h prior)
rescues VOR-increase learning specifically.} Left: rescue of the 1-Hz
VOR-increase learning deficit (KO vs.\ WT at 30 min post-diazepam
$p=0.86$); the rescue is absent 1 week later ($p=0.12$, 0 vs.\ 30 min).
Right: specificity. VOR-decrease learning and baseline VOR gain are
unchanged by the drug in either genotype (WT $p=0.91$; KO $p=0.37$;
KO vs.\ WT $p=0.11$; baseline gains $p\geq 0.36$).}
\label{fig:diazepam_rescue}
\end{figure}

\begin{figure}[htbp]
\centering
\includegraphics[width=0.95\textwidth]{figures/fig4_low_freq_no_rescue.pdf}
\caption{\textbf{Low-frequency (0.5 Hz) VOR-increase impairment is
\emph{not} rescued by behavioural pre-training or diazepam.} KO mice show
a trend toward impairment at 0.5 Hz ($p=0.06$ vs.\ WT). Neither VOR-decrease
pre-training ($p=0.47$), vestibular-only pre-training ($p=0.35$), nor
diazepam pre-treatment ($p=0.66$) reversed it.}
\label{fig:low_freq_no_rescue}
\end{figure}

\begin{figure}[htbp]
\centering
\includegraphics[width=0.95\textwidth]{figures/fig5_okr_rescue.pdf}
\caption{\textbf{OKR adaptation is rescued by diazepam, and post-drug KO
learning exceeds WT.} WT mice acquire the OKR adaptation over 60 min
($p=0.001$). KO mice fail post-saline ($p=0.87$, 0 vs.\ 60 min; KO vs.\ WT
$p=1.27\times10^{-4}$). 18--24~h post-diazepam, KO learning is rescued
(KO saline vs.\ diazepam $p=0.0001$) and exceeds WT (KO vs.\ WT $p=0.02$).
Baseline OKR gain is unaffected ($p=0.69$).}
\label{fig:okr_rescue}
\end{figure}

\begin{figure}[htbp]
\centering
\includegraphics[width=0.85\textwidth]{figures/fig6_metaplasticity_schematic.pdf}
\caption{\textbf{Threshold metaplasticity hypothesis.} In naive WT mice the
threshold for associative PF--Purkinje-cell LTD is normal, so spontaneous
circuit activity does not aberrantly induce LTD, and training recruits
LTD at the synapses where it supports learning. In L7-Fmr1 KO and
MHCI KbDb\textsuperscript{-/-} mice, enhanced LTD lowers the induction
threshold, allowing spontaneous activity to recruit LTD before training
and thereby raise the threshold for subsequent induction --- leaving LTD
unavailable to support learning during training. Pre-training or diazepam
suppresses activity-driven induction before training and resets the
threshold, restoring LTD-dependent learning.}
\label{fig:metaplasticity}
\end{figure}

\subsection{Behavioural pre-training rescues VOR-increase learning}

Having established that L7-Fmr1 KO mice recapitulate the task-selective
VOR phenotype of MHCI KbDb\textsuperscript{-/-} animals, we next asked
whether the same behavioural rescue strategy --- pre-training designed to
reverse PF--Purkinje cell LTD --- would also transfer to this molecularly
distinct genetic model. If the threshold metaplasticity hypothesis is
correct, then a manipulation that prevents or reverses aberrant LTD
induction during the period before training should restore the capacity
for subsequent LTD-dependent learning.
We implemented this manipulation behaviourally using VOR-decrease
pre-training, which is thought to engage PF--Purkinje cell LTP and thereby
reverse any pre-existing LTD \citep{levram2003reversing,
nguyenvu2017saturation}. In a within-mouse design, the same L7-Fmr1 KO
cohort that had failed to learn without pre-training was given 30 min of
VOR-decrease pre-training immediately before 30 min of VOR-increase
training.

The rescue was complete. In WT mice, both the pre-training and training
produced the expected adaptive changes (VOR decrease during pre-training,
$p = 0.02$, WT $-30$ vs. 0 min; VOR increase during training, $p = 0.001$,
WT 0 vs. 30 min). In L7-Fmr1 KO mice, the adaptive changes during both
periods were statistically indistinguishable from WT (VOR-decrease at the
end of pre-training, $p = 0.18$, KO vs. WT; VOR-increase learning from 0
to 30 min of training, $p = 0.17$, KO vs. WT, $t$-test). Crucially, the
same KO animals that had shown no change with VOR-increase training alone
($p = 0.99$, 0 vs. 30 min) now showed a significant increase when the
same training was delivered after VOR-decrease pre-training ($p = 0.001$,
0 vs. 30 min). The within-mouse comparison of learning without versus
with pre-training was itself significant ($p = 0.01$, paired-sample
$t$-test in KO animals) (Figure~\ref{fig:behavioral_pretraining}).

The intervention was specific to the KO phenotype. Pre-training did not
enhance WT learning beyond its already-significant baseline ($p = 0.41$,
paired $t$-test). A second behavioural pre-training protocol ---
vestibular-only stimulation in darkness, which habituates the VOR ---
produced an indistinguishable rescue (KO vs. WT at 30 min of VOR-increase
training, $p = 0.84$). To guard against a trivial carry-over from the
pre-training-induced decrease in the VOR, we subsampled the mice so that
the mean VOR-decrease at the end of pre-training was matched across
genotypes (within 2\%). The rescue was preserved (KO vs. WT, 30 min of
VOR-increase training: $p = 0.74$ after VOR-decrease pre-training;
$p = 0.40$ after vestibular-only pre-training), indicating that the KO
advantage comes from a genuine change in the learning capacity rather
than from a mean-level carry-over. Together, these findings demonstrate
that behavioural pre-training places the L7-Fmr1 KO cerebellum into a
state compatible with VOR-increase learning, just as it does in
MHCI KbDb\textsuperscript{-/-} animals \citep{nguyenvu2017saturation}.

\subsection{A single dose of diazepam rescues VOR-increase learning
18--24~h later}

The behavioural rescue demonstrates that the L7-Fmr1 KO cerebellum is not
intrinsically incapable of LTD-dependent learning --- the relevant
machinery remains intact and can be unlocked by an appropriate pre-training
protocol. However, behavioural pre-training reverses LTD in a task- and
region-specific manner --- a different pre-training protocol would be
required for each functional region of the cerebellum. We therefore asked whether systemic,
pharmacological suppression of neural activity --- which should limit
coincident parallel-fibre and climbing-fibre firing and thereby prevent
LTD induction throughout the cerebellum --- could achieve the same rescue.
Mice received a single IP injection of diazepam (0.5 mg/kg) or saline and
were tested 2~h, 18--24~h, or 1 week later.

The acute effect of diazepam was inhibitory, as expected. When VOR-increase
training was delivered 2~h after a 2.5 mg/kg or a 0.4--0.5 mg/kg dose,
learning was abolished in both WT and L7-Fmr1 KO mice (Figure~3 --- figure
supplement 2 in the experimental log). This result confirms pharmacological
action and is consistent with the requirement for normal neural activity
during training to induce LTD. Importantly, however, the baseline VOR in
the dark was unchanged at 2~h after diazepam (KO: $p = 0.72$, Pre vs. 2~h
Post; WT: $p = 0.99$), ruling out a gross effect on the reflex itself.

The diagnostic test is what happens after the drug has cleared. Eighteen
to twenty-four hours after the 0.5 mg/kg injection, the KO cohort that
had been unable to learn without the drug ($p = 0.99$, 0 vs. 30 min)
now showed robust VOR-increase learning indistinguishable from WT
($p = 0.86$, KO vs. WT at 30 min) (Figure~\ref{fig:diazepam_rescue}).
Baseline VOR gain remained unchanged at this time point (KO: $p = 0.77$,
Pre vs. 18--24~h post; WT: $p = 0.36$), so the rescue cannot be attributed
to a residual drug effect on the reflex itself. The rescue was transient:
one week after the same injection, the L7-Fmr1 KO mice again failed to
learn ($p = 0.12$, 0 vs. 30 min), indicating that diazepam restored a
functional state that does not persist and that must be re-triggered by
experience. WT mice did not show enhanced VOR-increase learning after
diazepam ($p = 0.55$ versus their own no-drug baseline; $p = 0.99$ across
sessions), consistent with a reset that is meaningful only for a circuit
in the aberrantly-recruited state. VOR-decrease learning was unchanged
by diazepam in either genotype (WT 1-day post-diazepam vs. control,
$p = 0.91$; KO, $p = 0.37$; KO vs. WT at 1-day post, $p = 0.11$), so the
drug's effect on learning was as selective as the deficit itself.

\subsection{The rescue is specific to LTD-dependent, high-frequency VOR
learning}

The diazepam rescue observed above could reflect a general pro-learning
effect unrelated to LTD. To probe the mechanism, we sought a case in which
the intervention should \emph{fail}: if diazepam operates through
LTD-specific metaplasticity, it should not rescue a VOR-increase deficit
that does not depend on PF--Purkinje cell LTD. Prior work shows that low-frequency ($\leq 0.66$~Hz)
VOR-increase learning is less dependent on LTD than its 1~Hz counterpart
\citep{boyden2004cerebellum}. We repeated the VOR-increase protocol at
0.5~Hz with both the behavioural and pharmacological rescues.

L7-Fmr1 KO mice showed a trend towards impairment at 0.5~Hz ($p = 0.06$,
KO vs. WT at 30 min), weaker than the high-frequency deficit. Neither
VOR-decrease pre-training (KO paired $t$-test, $p = 0.47$),
vestibular-only pre-training ($p = 0.35$), nor diazepam pre-treatment
($p = 0.66$, saline vs. diazepam, paired $t$-test) rescued this deficit
(Figure~\ref{fig:low_freq_no_rescue}). The frequency dissociation is a
double-check on the interpretation: the same interventions that reverse
the high-frequency deficit leave the low-frequency deficit untouched,
consistent with the high-frequency impairment being LTD-related and the
low-frequency component being the signature of an additional, LTD-
independent effect of \textit{Fmr1} loss in Purkinje cells.

\subsection{Diazepam rescues OKR adaptation and yields above-WT learning}

The frequency-selective dissociation from the previous subsection
strengthens the case for an LTD-specific mechanism. As a further, more
stringent test, we turned to OKR adaptation, the oculomotor task with the
strongest evidence for an obligate role of PF--Purkinje cell LTD \citep{inoshita2018okr,
takeuchi2008delphilin}, and therefore the strongest predicted rescue by
activity suppression. Sixty minutes of 1~Hz, $\pm 10^{\circ}/\text{s}$
optokinetic drum training induced robust learning in WT animals
($p = 0.001$, 0 vs. 60 min). In L7-Fmr1 KO animals after saline, OKR
adaptation was essentially abolished ($p = 0.87$, 0 vs. 60 min;
$p = 1.27 \times 10^{-4}$ for KO vs. WT at 60 min). Eighteen to twenty-four
hours after a single 0.5 mg/kg diazepam injection, the same KO cohort
showed complete rescue (KO saline vs. diazepam pre-treatment at 60 min,
$p = 0.0001$), and notably, their post-diazepam learning exceeded WT
(KO vs. WT at 60 min, post-diazepam, $p = 0.02$)
(Figure~\ref{fig:okr_rescue}). Baseline OKR gain was unaffected by
genotype (Pre, $p = 0.69$) or by diazepam (KO post-saline vs.
post-diazepam, $p = 0.69$; WT, $p = 0.55$).

The above-WT learning in KO animals after diazepam is a striking result.
It suggests that the enhanced LTD which has been aberrantly recruited in
the untreated KO cerebellum can, once reset by activity suppression, be
redirected to support learning more effectively than the unperturbed WT
machinery. This is consistent with the idea that the KO cerebellum retains
--- and in fact has more than the usual amount of --- LTD capacity, but
that in ordinary conditions that capacity is not available for training-
driven plasticity. The schematic in
Figure~\ref{fig:metaplasticity} summarises the hypothesis: enhanced
LTD lowers the induction threshold and allows spontaneous activity to
aberrantly engage the plasticity machinery; pre-training or diazepam
resets the threshold, making LTD available to support learning during
training.

\begin{table}[htbp]
\centering
\caption{Summary of extracted oculomotor-learning statistics. Unless noted,
$p$-values are from Tukey post-hoc tests following two-factor
repeated-measures ANOVA. Reported to the precision in the experimental log.}
\label{tab:summary}
\small
\begin{tabular}{p{2.6cm}cccp{2.3cm}c}
\toprule
Task & Stim. & Condition & Genotype & Comparison & $p$ \\
\midrule
VOR-increase & 1 Hz & No drug & WT  & 0 vs 30 min & $7.37{\times}10^{-4}$ \\
VOR-increase & 1 Hz & No drug & KO  & 0 vs 30 min & $0.97$ \\
VOR-increase & 1 Hz & No drug & KO vs WT & 30 min   & $3.47{\times}10^{-5}$ \\
VOR-decrease & 1 Hz & No drug & WT  & 0 vs 30 min & $0.001$ \\
VOR-decrease & 1 Hz & No drug & KO  & 0 vs 30 min & $1.10{\times}10^{-5}$ \\
VOR-decrease & 1 Hz & No drug & KO vs WT & 30 min   & $0.091$ \\
VOR-increase & 1 Hz & +VOR-dec pre-train & KO  & 0 vs 30 min & $0.001$ \\
VOR-increase & 1 Hz & +Vest-only pre-train & KO vs WT & 30 min & $0.84$ \\
VOR-increase & 1 Hz & KO paired & KO & pre vs no-pre & $0.01$ \\
VOR-increase & 1 Hz & 18--24 h post-diazepam & KO vs WT & 30 min & $0.86$ \\
VOR-increase & 1 Hz & 1 week post-diazepam & KO  & 0 vs 30 min & $0.12$ \\
VOR-decrease & 1 Hz & 1 d post-diazepam & KO paired & drug vs ctrl & $0.37$ \\
VOR-decrease & 1 Hz & 1 d post-diazepam & KO vs WT & 30 min & $0.11$ \\
VOR-increase & 0.5 Hz & No drug & KO vs WT & 30 min & $0.06$ \\
VOR-increase & 0.5 Hz & +diazepam & KO paired & saline vs drug & $0.66$ \\
OKR         & 1 Hz & Saline & WT  & 0 vs 60 min & $0.001$ \\
OKR         & 1 Hz & Saline & KO  & 0 vs 60 min & $0.87$ \\
OKR         & 1 Hz & Saline & KO vs WT & 60 min & $1.27{\times}10^{-4}$ \\
OKR         & 1 Hz & +diazepam & KO paired & saline vs drug & $0.0001$ \\
OKR         & 1 Hz & +diazepam & KO vs WT & 60 min & $0.02$ \\
Baseline OKR & 1 Hz & Pre-training & KO vs WT & gain & $0.690$ \\
\bottomrule
\end{tabular}
\end{table}

\section{Discussion}

The central finding of this study is that enhancement of associative
plasticity in a second, molecularly distinct genetic model reproduces the
same paradoxical learning impairment previously described in
MHCI KbDb\textsuperscript{-/-} animals, and that this impairment can be
eliminated by a single dose of a widely used GABAergic drug given well
before training. L7-Fmr1 KO mice were selectively impaired on 1~Hz
VOR-increase learning ($p = 0.97$, 0 vs. 30 min) and OKR adaptation
($p = 1.27 \times 10^{-4}$ vs. WT at 60 min), with intact VOR-decrease
learning. Both deficits were completely reversed by 18--24~h
pre-treatment with 0.5 mg/kg diazepam, and the OKR rescue went even
further: post-diazepam learning in KO mice exceeded WT ($p = 0.02$).
These patterns cannot be parsimoniously explained by acute pharmacology
(the drug had cleared by the time of training), by gross motor or
vestibular effects (baseline VOR and OKR gain were unchanged), or by
generic enhancement of plasticity (WT learning was unaffected by diazepam).
They converge instead on the account proposed by
\citet{nguyenvu2017saturation}, in which the enhancement of LTD makes it
aberrantly recruitable by ongoing activity, elevating the threshold for
its training-driven induction and thereby making it unavailable to support
learning --- a state that can be reset by activity suppression during a
pre-training window.

The convergence across models with entirely different molecular origins
is the strongest evidence that the phenotype tracks enhanced LTD per se
rather than some coincident consequence of the specific genetic lesion.
MHCI H2-K\textsuperscript{b} and H2-D\textsuperscript{b} act through
MAP-kinase and integrin pathways via immune receptors such as PirB,
whereas FMRP acts as an RNA-binding translational regulator whose loss
derepresses mGluR-dependent dendritic protein synthesis
\citep{koekkoek2005fmr1, nosyreva2006fmr1}. That two such distinct
molecular perturbations produce the same task-selective oculomotor deficit
--- and the same rescue by two unrelated interventions (behavioural
pre-training and diazepam) --- is difficult to reconcile with an off-target
explanation. It is, however, exactly what the threshold metaplasticity
hypothesis predicts: if the shared downstream feature is enhanced
PF--Purkinje cell LTD, then the shared behavioural consequence of its
interaction with ongoing circuit activity, and the shared efficacy of
interventions that limit that interaction, follows naturally.

Task-specificity within the oculomotor repertoire is an equally stringent
constraint. The cerebellar flocculus processes the same vestibular,
visual, and motor signals for VOR-increase, VOR-decrease, and OKR learning,
so a non-specific impairment would not distinguish among them
\citep{boyden2004cerebellum, raymond2018cerebellum}. The learning that
was impaired in the KO and rescued by our interventions (high-frequency
VOR-increase; OKR adaptation) is exactly the set that prior LTD-focused
work has identified as LTD-dependent
\citep{inoshita2018okr, takeuchi2008delphilin, kimpo2014gating,
rowan2018graded}, while VOR-decrease learning and low-frequency (0.5 Hz)
VOR-increase learning --- tasks that are less dependent on LTD
\citep{kimpo2005vor} --- were not rescued. This within-animal dissociation
also cautions against a reading of the results that attributes Fragile X
cerebellar learning deficits solely to enhanced LTD: the persistent
low-frequency impairment that survives our interventions (KO vs. WT at
0.5 Hz, $p = 0.06$; diazepam rescue, $p = 0.66$) likely reflects an
additional Purkinje-cell phenotype that \textit{Fmr1} loss produces on
top of LTD enhancement. Characterising this residual phenotype is an
important direction for future work, particularly given the clinical
interest in cerebellar contributions to FXS
\citep{koekkoek2005fmr1}.

The clinical implication is encouraging but should be read with care.
Diazepam at 0.5 mg/kg --- a low dose relative to clinical use --- produced
complete behavioural rescue without measurable side effects on baseline
oculomotor performance or on VOR-decrease learning, which likely relies on
PF--Purkinje cell LTP at the same synapses where LTD supports VOR-increase
learning \citep{levram2003reversing}. The selectivity of the intervention
at the behavioural level matters for repurposing. Because pre-training
activity suppression should limit LTD induction wherever LTD operates in
the cerebellum, the strategy generalises in principle to cerebellar learning
tasks beyond the oculomotor domain, and possibly to other brain regions in
which enhanced associative LTD has been implicated
\citep{raymond2018cerebellum}. The transient nature of the rescue (absent
one week later) tempers expectations: a clinical strategy would require a
scheduled, pre-training dose, not a one-time cure, and would need to be
balanced against the well-known cognitive side effects of chronic
benzodiazepine use. Nevertheless, the demonstration that a single, low,
subacute dose of an approved drug can restore a genetic learning deficit
in a mouse model motivates serious consideration of this class of
interventions for disorders in which enhanced associative plasticity is
thought to contribute to learning impairment.

The theoretical upshot bears on the BCM framework. In its classical form,
a sliding threshold for Hebbian LTP provides a homeostatic counterweight
to runaway potentiation and firing rates \citep{abraham1996metaplasticity}.
The plasticity at PF--Purkinje cell synapses is anti-Hebbian: associative
LTD depresses synapses whose activity coincides with a climbing-fibre
error signal, and a sliding threshold for LTD therefore cannot serve
firing-rate homeostasis in the usual sense --- a lower threshold for LTD
in response to higher activity would actively destabilise Purkinje-cell
firing rates. Our results, in combination with prior work
\citep{nguyenvu2017saturation}, imply that threshold metaplasticity at
these synapses nevertheless operates, but in the service of a different
function: limiting plasticity once it has been recruited, thereby
protecting recent memories from being overwritten and potentially
separating memories acquired in close succession onto distinct Purkinje
cells, in partial analogy with activity-history-dependent memory-
allocation mechanisms described in other circuits \citep{sargin2013creb}.
The observation that the metaplastic state is driven by the recent history
of plasticity induction, not by firing rate per se --- evidenced by the
different effects of pre-training on WT and KO animals --- is consistent
with candidate substrates that include Purkinje-cell dendritic excitability
\citep{ohtsuki2012sk2}, molecular-layer interneuron inhibition
\citep{rowan2018graded}, and AMPA-receptor trafficking via stargazin
dephosphorylation \citep{nomura2012stargazin}. Testing these substrates
for activity-dependent modulation in L7-Fmr1 KO animals is a natural next
step.

A limitation of the present work is that we do not directly measure the
threshold for LTD induction at PF--Purkinje cell synapses during the
pre-training window. The threshold metaplasticity hypothesis is supported
here by the specificity of its behavioural predictions --- by molecule,
by task, by frequency, by time course --- but a physiological measurement
of threshold changes during and after pre-training or diazepam is needed
for a full test. It will also be important to extend the pharmacological
strategy to cerebellar tasks beyond the oculomotor domain, to other
regions with pathologically enhanced LTD, and to chronic rather than
acute intervention schedules. More broadly, our results highlight a
principle that may generalise beyond cerebellum: synaptic plasticity and
the pattern of ongoing circuit activity together determine the capacity
for new learning, so the restoration of learning after a plasticity-
enhancing lesion need not require correcting the lesion itself, but may
instead be achieved by appropriately shaping the circuit's recent activity
history. Learning deficits associated with enhanced plasticity need not
be permanent.

\bibliographystyle{unsrtnat}
\bibliography{references}

\end{document}
